% Doku — LuaLaTeX preamble injected after pandoc's own preamble.
% The typographic profile is the printed twin of the preview page:
%   - type scale 2.2 / 1.55 / 1.25 / 1.1 of an 11pt body, rule under H2,
%     line height 1.65 (linestretch set by the exporter), paper/ink palette,
%     body text in the preview's secondary ink, headings in the primary ink;
%   - block quotes on a tinted card with an accent bar (mdframed), code on
%     the warm code shade (fancyvrb frame), inline code highlighted;
%   - tables with a shaded bold header row and hairline rules, sized by the
%     shared column algorithm (tableWidths.lua);
%   - task-list boxes drawn like the preview's, hairline horizontal rules;
%   - nothing is ever clipped: code lines and long tokens wrap.
% The exporter declares A4 twoside with the preview's binding margins.

\usepackage{microtype}
% No character protrusion: every glyph stays inside the binding-safe text
% band, exactly like the preview page.
\microtypesetup{protrusion=false}
\usepackage{etoolbox}
\usepackage{calc}
\usepackage{longtable,booktabs,array}
\usepackage{colortbl}
\usepackage{fvextra}
\usepackage{luacolor}
\usepackage{lua-ul}
\usepackage[framemethod=tikz]{mdframed}
\usepackage{titlesec}
\usepackage{fancyhdr}
\usepackage{setspace}

% ------------------------------------------------------------- palette
\definecolor{dokuaccent}{HTML}{00A3EE}
\definecolor{dokulink}{HTML}{007EB8}
\definecolor{dokuink}{HTML}{1A1816}
\definecolor{dokutext}{HTML}{5F5953}
\definecolor{dokurule}{HTML}{D9D3CA}
\definecolor{dokushade}{HTML}{F4EFE8}
\definecolor{dokuinline}{HTML}{EEF8FF}
\definecolor{dokuquote}{HTML}{F9FCFE}
\definecolor{dokumuted}{HTML}{8A837B}
% Pandoc's own code background (framed's snugshade) uses shadecolor.
\definecolor{shadecolor}{HTML}{F4EFE8}

% Body text in the preview's secondary ink; headings restore the primary ink.
\AtBeginDocument{\color{dokutext}}

% ------------------------------------------------------------ headings
% Headings are set ragged right like the preview, so a long word breaks at
% its discretionary points instead of running into the margin. The exporter
% sets linestretch 1.375 (the preview's 1.65 line height for the body); the
% leading below is pre-divided by that stretch so headings keep the
% preview's 1.15 leading without touching \baselinestretch (changing it inside
% a title format misplaces the block that follows under fontspec).
\titleformat{\section}{\fontsize{24.2pt}{20.4pt}\selectfont\bfseries\color{dokuink}\raggedright}{}{0pt}{}
\titleformat{\subsection}{\fontsize{17pt}{15.3pt}\selectfont\bfseries\color{dokuink}\raggedright}{}{0pt}{}[{\color{dokurule}\titlerule}]
\titleformat{\subsubsection}{\fontsize{13.75pt}{12.4pt}\selectfont\bfseries\color{dokuink}\raggedright}{}{0pt}{}
\titleformat{\paragraph}[hang]{\fontsize{12.1pt}{10.9pt}\selectfont\bfseries\color{dokuink}\raggedright}{}{0pt}{}
\titleformat{\subparagraph}[hang]{\fontsize{11pt}{10pt}\selectfont\bfseries\color{dokuink}\raggedright}{}{0pt}{}
\titlespacing*{\section}{0pt}{1.6em}{0.5em}
\titlespacing*{\subsection}{0pt}{1.6em}{0.5em}
\titlespacing*{\subsubsection}{0pt}{1.6em}{0.5em}
\titlespacing*{\paragraph}{0pt}{1.4em}{0.4em}
\titlespacing*{\subparagraph}{0pt}{1.2em}{0.3em}

% Paragraphs: no indent, one line of air between blocks, like the preview.
\setlength{\parindent}{0pt}
\setlength{\parskip}{0.85em plus 0.15em minus 0.1em}
% Pathological tokens (hashes, URLs, joined words) get discretionary break
% points from tableWidths.lua; let the third pass loosen a line rather than
% push text into the margin.
\setlength{\emergencystretch}{6em}

% ---------------------------------------------------------------- page
\pagestyle{fancy}
\fancyhf{}
\renewcommand{\headrulewidth}{0pt}
\fancyfoot[C]{\small\color{dokumuted}\thepage}
\fancypagestyle{plain}{\fancyhf{}\renewcommand{\headrulewidth}{0pt}\fancyfoot[C]{\small\color{dokumuted}\thepage}}

% --------------------------------------------------------------- lists
\renewcommand{\labelitemii}{$\circ$}
\renewcommand{\labelitemiii}{\textbullet}

% Task-list boxes: pandoc writes $\square$ / $\boxtimes$; draw them like the
% preview's check boxes (hairline box, filled box with a white tick).
\newcommand{\dokubox}{\tikz[baseline=-0.05em]{\draw[dokurule,line width=0.6pt] (0,0) rectangle (0.85em,0.85em);}}
\newcommand{\dokuboxchecked}{\tikz[baseline=-0.05em]{\fill[dokulink] (0,0) rectangle (0.85em,0.85em);\draw[white,line width=1pt] (0.2em,0.42em) -- (0.36em,0.24em) -- (0.68em,0.64em);}}
\AtBeginDocument{%
  \renewcommand{\square}{\text{\dokubox}}%
  \renewcommand{\boxtimes}{\text{\dokuboxchecked}}%
}

% Horizontal rule: full-width hairline in the rule colour (via tableWidths.lua).
\newcommand{\dokurule}{\par\vspace{0.4em}{\color{dokurule}\noindent\rule{\linewidth}{0.5pt}}\par\vspace{0.4em}}

% ---------------------------------------------------------- block quote
% Tinted card with the accent bar on the left, splittable across pages.
% Every mdframed style lists all of its line options: an option set by one
% frame otherwise leaks into the next one (a quote's `topline=false` would
% shift the following code box up over its heading).
\mdfdefinestyle{dokuquote}{backgroundcolor=dokuquote,linecolor=dokuaccent,linewidth=2.5pt,
  leftline=true,rightline=false,topline=false,bottomline=false,roundcorner=0pt,
  innerleftmargin=0.9em,innerrightmargin=1em,innertopmargin=0.7em,innerbottommargin=0.7em,
  skipabove=0.6em,skipbelow=1em,leftmargin=0pt,rightmargin=0pt,nobreak=false}
\renewenvironment{quote}{%
  \par\vspace{0.35em}\begin{mdframed}[style=dokuquote]\color{dokutext}\ignorespaces}%
  {\end{mdframed}}

% ---------------------------------------------------------------- code
% Plain verbatim (no syntax colours, like the preview) in the code font at
% 0.95em, every line wrapped, on the warm shade with rounded corners.
% Code at 0.95em with the preview's 1.5 line height (15.7pt after the body
% stretch), on
% the warm shade drawn by fancyvrb itself (frame and fill in the same colour):
% a frame package around verbatim (mdframed, tcolorbox) misplaces the box
% over the heading above it, fancyvrb's own frame never does and splits
% across pages line by line.
\newcommand{\dokucodesize}{\fontsize{10.45pt}{11.42pt}\selectfont}
\newcommand{\dokucodeframe}{frame=single,framerule=0.4pt,framesep=0.8em,rulecolor=\color{dokushade},fillcolor=\color{dokushade}}
\fvset{breaklines=true,breakanywhere=true,breakindent=0pt,breaksymbolleft={},breaksymbolright={},fontsize=\dokucodesize,formatcom=\color{dokutext}}
\RecustomVerbatimEnvironment{verbatim}{Verbatim}{breaklines=true,breakanywhere=true,breakindent=0pt,breaksymbolleft={},breaksymbolright={},fontsize=\dokucodesize,formatcom=\color{dokutext},frame=single,framerule=0.4pt,framesep=0.8em,rulecolor=\color{dokushade},fillcolor=\color{dokushade}}
\makeatletter
\@ifundefined{Highlighting}{}{%
  \DefineVerbatimEnvironment{Highlighting}{Verbatim}{commandchars=\\\{\},breaklines,breakanywhere,breakindent=0pt,breaksymbolleft={},breaksymbolright={},fontsize=\dokucodesize,frame=single,framerule=0.4pt,framesep=0.8em,rulecolor=\color{dokushade},fillcolor=\color{dokushade}}}
\makeatother
% Air around the box, like the preview's margins.
\BeforeBeginEnvironment{verbatim}{\par\vspace{0.2em}}
\AfterEndEnvironment{verbatim}{\vspace{0.4em}}

% Inline code: highlighted pill that can still break across lines.
\let\dokuplaintt\texttt
\renewcommand{\texttt}[1]{\highLight[dokuinline]{\dokuplaintt{#1}}}

% -------------------------------------------------------------- tables
% Table text at 0.95em, hairline rules in the paper palette, shaded header
% row (the first rule of every head is followed by a row colour; pandoc emits
% \toprule at the start of the first and of every repeated head).
% Table text at 0.95em with the preview's 1.5 line height (15.7pt once the
% body stretch of 1.375 is applied).
\newcommand{\dokutablesize}{\fontsize{10.45pt}{11.42pt}\selectfont}
\arrayrulecolor{dokurule}
\setlength{\heavyrulewidth}{0.5pt}
\setlength{\lightrulewidth}{0.5pt}
\setlength{\abovetopsep}{0pt}
\setlength{\belowbottomsep}{0pt}
\renewcommand{\arraystretch}{1.3}
\setlength{\LTpre}{0.6em}
\setlength{\LTpost}{1em}
\AtBeginEnvironment{longtable}{\dokutablesize}
\let\dokutoprule\toprule
\renewcommand{\toprule}{\dokutoprule\rowcolor{dokushade}}

% Pages are not vertically justified: like the preview, blocks keep their
% own spacing and the last page line simply ends where the content ends.
\raggedbottom
